\section*{Introduction}

\textbf{Objectif :} 
Construction d'un modèle 1D homogénéisé de transport
et réaction dans une électrode poreuse (\ce{CO2RR}), à partir d'un modèle
3D simple à l'échelle du pore. Des possibilités pour améliorer le modèle
 de base sont également discutées. 
 %Est ce qu'on dit qu'on vut un modèle qui donne des résultats cohérents 
 %avec les expériences mais sans truc qui n'ajoute rien
 
\textbf{Contenu :}
\begin{itemize}
    \item Modèle 3D local : Nernst--Planck (espèces), Ohm (électrons),
    Butler--Volmer (cinétique), électroneutralité / Poisson (potentiel).
    \item Réduction 1D par homogénéisation (prise de moyenne volumique,
    coefficients effectifs), dont le modèle de Thiele est obtenu comme cas
    limite (solution analytique, facteur d'efficacité).
    \item Cas test : diffusion--réaction \ce{CO2} $\rightarrow$ \ce{CO}.
    \item Limites et extensions du modèle de base.
    \item Nombres adimensionnels : Péclet, Damköhler, Debye.
\end{itemize}

\textbf{Annexes :}
\begin{itemize}
  \item Rappels théoriques : moyenne superficielle, théorème de la divergence de Gauss, théorème de la moyenne spatiale, théorème de transport.
  \item Liste des grandeurs et paramètres utilisés dans le modèle.
\end{itemize}

%Je pourrais potentielement faire une vraie intro + une table des matièères
%mais je ne suis pas sûre que ce soit pertinent ic
%En effet, le but de ce document est de présenter le modèle et les résultats, pas de faire un mémoire complet sur le sujet.